\SourceLeaf{434}%
hence
\[
\begin{aligned}
\frac{\partial X}{\partial x}&=Z^{3}\{-(1+uu)T+tuU\},&\quad\frac{\partial Y}{\partial x}&=Z^{3}\{tuT-(1+tt)U\},\\
\frac{\partial X}{\partial y}&=Z^{3}\{-(1+uu)U+tuV\},&\frac{\partial Y}{\partial y}&=Z^{3}\{tuU-(1+tt)V\}
\end{aligned}
\]
and
\[
\begin{aligned}
\omega&=Z^{6}(TV-UU)((1+tt)(1+uu)-ttuu)\\
&=Z^{6}(TV-UU)(1+tt+uu)\\
&=Z^{4}(TV-UU)\\
&=\frac{TV-UU}{(1+tt+uu)^{2}},
\end{aligned}
\]
exactly the same expression that we found at the end of the preceding article. We see, therefore, that
\begin{quote}
»the measure of curvature is always expressed by the fraction whose numerator is $=1$ and whose denominator is the product of the greatest and least radii of curvature in the planes passing through the normal«.
\end{quote}

\Needspace{5\baselineskip}
\begin{center}12.\end{center}
\noindent\hspace*{1.2em}We shall now investigate the nature of the shortest line on the curved surface. The nature of a curved line in space is determined in general by regarding the coordinates $x,y,z$ of each point as functions of a single variable, which we shall call $w$. The length of the curve up to this point, measured from an arbitrary starting point, is then
\[
=\int\BareRoot{\left\{\left(\frac{dx}{dw}\right)^{2}+\left(\frac{dy}{dw}\right)^{2}+\left(\frac{dz}{dw}\right)^{2}\right\}}.dw.
\]
If the position of the curve is changed by an infinitesimal variation, the variation of its total length is
